% =========================================================
% 1.3 OBJETIVOS
% El \section{} de esta sección está en el chapter.tex del capítulo;
% sus referencias van en seccion1_3.bib (misma carpeta).
% =========================================================

\subsection{Objetivo general}
\label{subsec:objetivo_general}

Desarrollar un prototipo de aplicación móvil que genere rutas óptimas para peatones en una zona delimitada de la alcaldía Gustavo A. Madero mediante el uso de un algoritmo de enrutamiento multicriterio basado en factores de diseño ambiental y dinámicas temporales para minimizar la exposición al riesgo. 

\subsection{Objetivos específicos}
\label{subsec:objetivos_especificos}

\begin{enumerate}
    \item Modelar la red vial de la alcaldía Gustavo A. Madero como un grafo ponderado, integrando datos geoespaciales y delictivos, para establecer la base de datos del sistema.
    \item Implementar un algoritmo de enrutamiento multicriterio, evaluando dinámicamente los pesos del grafo, para generar y clasificar alternativas de rutas seguras.
    \item Desarrollar la interfaz gráfica del prototipo de la aplicación móvil, utilizando servicios de mapas y geolocalización, para permitir la consulta de trayectos y la captura de retroalimentación del usuario.
    \item Evaluar el desempeño del prototipo mediante pruebas, para validar su efectividad en la evasión de zonas que propician el delito.
\end{enumerate}
